(1) 在球内 \(+z\) 轴设置镜像电荷 \(q'\) 于 \(a=R^2/d<R\)。对球面任意极角，点到真实电荷距离 \(s=\sqrt{R^2+d^2-2Rd\cos\theta}\)，到镜像位置距离 \(s'=\sqrt{R^2+a^2-2Ra\cos\theta}=(R/d)s\)。球面要求 \(q/s+q'/s'=0\)，因此 \(q'=-qR/d\)。导体外的 Dirichlet 边界问题解唯一，所以该构造给出完整外部静电场，而非仅在轴上的近似。

(2) 导体外任意位置 \(\mathbf r\) 的电势为 \(\varphi(\mathbf r)=\frac{1}{4\pi\varepsilon_0}[q/|\mathbf r-d\hat z|-(qR/d)/|\mathbf r-(R^2/d)\hat z|]\)。在球面利用 \(s'=(R/d)s\) 两项逐点抵消，\(\varphi=0\)。当观测距离 \(r\gg d\) 时，主导远场电势为 \([q+q']/ (4\pi\varepsilon_0r)=q(1-R/d)/(4\pi\varepsilon_0r)\)，由真实点电荷与球表面总感应电荷构成整体有效单极矩。

(3) 镜像法在导体球外给出与真实感应电荷相同的电场，因此点电荷受力等效于镜像电荷作用，沿球心方向为负。轴向距离为 \(d-a=(d^2-R^2)/d\)，故 \(F_z=\frac1{4\pi\varepsilon_0}qq'/(d-a)^2=-q^2Rd/[4\pi\varepsilon_0(d^2-R^2)^2]\)。由于镜像电荷只是边界问题的数学构造，计算真实系统能量时不能把两电荷的普通两体势能简单照搬而不考虑感应源随 \(d\) 的变化；力可直接由镜像场获得。

(4) 对导体球面取紧邻外侧的高斯面，外部电场通量与球内镜像源等价，故真实球的感应总电荷 \(Q_{\rm ind}=q'=-qR/d<0\)，符合接地导体从地面汲取负电荷的物理图景。若 \(d\to\infty\)，\(Q_{\rm ind}\to0\)、\(F_z\sim-q^2R/(4\pi\varepsilon_0d^3)\)；若 \(d\to R^+\)，感应总电荷趋于 \(-q\)，吸引力在点电荷离表面间距趋零时发散，说明点电荷与理想连续导体表面近距离模型需受微观尺度修正。
